\documentclass[11pt,letterpaper]{article}

\usepackage[spanish,es-noquoting]{babel}
\usepackage[utf8]{inputenc}
\usepackage[T1]{fontenc}
\usepackage{lmodern}
\usepackage{amsmath,amssymb,mathtools}
\usepackage{amsthm}
\usepackage{physics}
\usepackage{geometry}
\usepackage{enumitem}
\usepackage{needspace}
\usepackage{tikz}
\usepackage{pgfplots}
\pgfplotsset{compat=1.18}
\definecolor{diagramblue}{RGB}{0,142,190}
\definecolor{diagramgreen}{RGB}{67,150,45}
\definecolor{diagramred}{RGB}{210,40,30}
\usetikzlibrary{arrows.meta,calc,positioning}
\usepackage{hyperref}

\hypersetup{
  pdftitle={Teoría lambda phi4 y reglas de Feynman en espacio de posiciones},
  pdfauthor={Omar Corona Tejeda},
  colorlinks=true,
  linkcolor=blue,
  urlcolor=blue,
  citecolor=blue
}

\geometry{margin=2.2cm}
\setlength{\parindent}{0pt}
\setlength{\parskip}{0.45em}
\interfootnotelinepenalty=10000
\raggedbottom
\clubpenalty=10000
\widowpenalty=10000
\setlist[itemize]{leftmargin=1.3em}

\theoremstyle{definition}
\newtheorem{definicion}{Definición}
\theoremstyle{definition}
\newtheorem{observacion}{Observación}
\newtheorem{ejemplo}{Ejemplo}
\AddToHook{env/ejemplo/begin}{\pushQED{\qed}}
\AddToHook{env/ejemplo/end}{\par\nobreak\vspace{0.35em}\noindent\popQED}
\newcommand{\eqbox}[1]{\boxed{\displaystyle #1}}
\newcounter{footnotequation}
\newcommand{\fnEqTag}{\refstepcounter{footnotequation}\tag{\Roman{footnotequation}}}
\newcounter{diagrama}
\newcommand{\diagramcaption}[1]{%
  \refstepcounter{diagrama}%
  \par\vspace{0.7em}{\footnotesize\textbf{Figura \thediagrama.} #1}\par\vspace{0.4em}%
}
\newcommand{\qftFreeDiagram}{%
  \begin{tikzpicture}[baseline=-0.5ex,scale=0.45]
    \draw[blue,very thick] (-1,0)--(1,0);
  \end{tikzpicture}%
}
\newcommand{\qftTadpoleDiagram}{%
  \begin{tikzpicture}[baseline=-0.5ex,scale=0.45]
    \draw[blue,very thick] (-1,0)--(1,0);
    \draw[blue,very thick] (0,0)
      .. controls (-0.6,1.1) and (0.6,1.1) .. (0,0);
    \fill[red] (0,0) circle (2pt);
  \end{tikzpicture}%
}
\newcommand{\qftVacuumDiagram}{%
  \begin{tikzpicture}[baseline=-0.5ex,scale=0.35]
    \draw[blue,very thick] (0,0)
      .. controls (-0.7,1.3) and (0.7,1.3) .. (0,0);
    \draw[blue,very thick] (0,0)
      .. controls (-0.7,-1.3) and (0.7,-1.3) .. (0,0);
    \fill[red] (0,0) circle (2pt);
  \end{tikzpicture}%
}
\newcommand{\qftChainDiagram}{%
  \begin{tikzpicture}[baseline=-0.5ex,scale=0.45]
    \draw[blue,very thick] (-1.4,0)--(1.4,0);
    \foreach \x in {-0.5,0.5} {
      \draw[blue,very thick] (\x,0)
        .. controls (\x-0.5,1.1) and (\x+0.5,1.1) .. (\x,0);
      \fill[red] (\x,0) circle (2pt);
    }
  \end{tikzpicture}%
}

\title{Teoría $\lambda\phi^4$ y reglas de Feynman\\
en espacio de posiciones}
\author{Omar Corona Tejeda\\
\href{mailto:omarct1989@ciencias.unam.mx}{omarct1989@ciencias.unam.mx}\\[0.2cm]
\href{https://artinnether.github.io/the_mathematical_physics_project/}
{The Mathematical Physics Project}}
\date{30 de septiembre de 2026}

\begin{document}

\maketitle

\begin{abstract}
Se desarrolla una derivación de las reglas de Feynman en espacio de
posiciones para el correlador de dos puntos de la teoría escalar real
$\lambda\phi^4$. Primero se precisan las imágenes de Schrödinger, Heisenberg
e interacción, distinguiendo la evolución libre de la evolución completa.
Mediante conmutadores anidados y la fórmula de Baker--Campbell--Hausdorff,
se reconstruye la expansión modal del campo libre y se obtiene el
Hamiltoniano de interacción.

La fórmula de Gell-Mann--Low, la serie de Dyson y el teorema de Wick
permiten traducir las contracciones en propagadores, vértices e integrales
sobre el espacio-tiempo. Se explica la cancelación de las burbujas de vacío,
se calculan las contribuciones de primer orden y se presentan las tres
topologías conectadas de segundo orden con sus integrales y factores de
simetría. Las figuras y los ejemplos conectan la evolución del campo con el cálculo
operatorial y muestran cómo las reglas diagramáticas surgen de la
combinatoria de contracciones. Las integrales de lazos se
mantienen formales, sin desarrollar su regularización y renormalización.
\end{abstract}

\section{Introducción}

El paso de la dinámica de un campo cuántico a sus reglas de Feynman exige
articular varias estructuras: la evolución de estados y operadores, la
elección del vacío, el ordenamiento temporal y la combinatoria de las
contracciones. La sencillez de un diagrama es el resultado de ese trabajo.
Comprender cómo se construye permite reconocer qué representa cada línea,
qué operación introduce un vértice y de dónde proviene el coeficiente que
acompaña a una contribución. Estas notas desarrollan esa construcción para
el correlador de dos puntos de la teoría escalar real $\lambda\phi^4$ en
espacio de posiciones.

El primer problema es identificar los operadores con los que se calcula.
Las imágenes de Schrödinger, Heisenberg e interacción comparten los datos
iniciales, pero distribuyen de manera distinta la evolución temporal entre
estados y operadores. Distinguir el Hamiltoniano libre del Hamiltoniano
completo permite establecer el papel de $\phi_I$ y su relación con el campo
interactuante $\phi_H$. La expansión modal se reconstruye mediante
conmutadores anidados y la fórmula de Baker--Campbell--Hausdorff; así se
verifica explícitamente la evolución libre que sustenta el desarrollo
perturbativo y se construye el Hamiltoniano de interacción.

El segundo problema es convertir esa estructura operatorial en un cálculo
organizado. La fórmula de Gell-Mann--Low relaciona el correlador del vacío
interactuante con valores esperados en el vacío libre; la serie de Dyson
ordena las inserciones de la interacción y el teorema de Wick las reduce
a productos de propagadores. Seguimos las contracciones hasta obtener las
contribuciones de primer orden y las tres topologías conectadas de segundo
orden, con sus integrales y factores de simetría. La normalización permite
establecer qué piezas se cancelan: las burbujas de vacío desaparecen,
mientras que los lazos ligados a los puntos externos permanecen.

El resultado central es una justificación de las reglas de Feynman desde
el formalismo canónico y el conteo de contracciones. Una vez establecida
esta base, los cálculos posteriores pueden organizarse mediante diagramas,
sin reconstruir en cada caso toda la expansión operatorial. Trabajamos con
$\hbar=c=1$, métrica $(+,-,-,-)$ y una interacción sin orden normal,
alrededor del vacío simétrico. Las integrales de lazos se mantienen formales;
su evaluación, regularización y renormalización constituyen la siguiente
etapa del análisis.

\section{Teoría \texorpdfstring{$\lambda\phi^4$}{lambda phi4}}

Consideramos un campo escalar real en el régimen perturbativo alrededor
del vacío simétrico, con $m^2>0$ y $\lambda\geq0$. Su lagrangiano es
\begin{equation}
  \eqbox{\mathcal{L}
  \coloneqq\frac{1}{2}\partial_\mu\phi\,\partial^\mu\phi
  -\frac{m^2}{2}\phi^2
  -\frac{\lambda}{4!}\phi^4.}
\end{equation}
El punto sobre el campo denota su derivada parcial respecto del tiempo,
manteniendo fija la coordenada espacial:
\begin{equation}
  \dot\phi(t,\mathbf{x})
  \coloneqq\partial_t\phi(t,\mathbf{x})
  =\frac{\partial\phi(t,\mathbf{x})}{\partial t}
  =\partial_0\phi(t,\mathbf{x}).
\end{equation}
Con la convención métrica $(+,-,-,-)$, el momento canónico conjugado al campo es
\begin{equation}
  \pi(t,\mathbf{x})
  \coloneqq\frac{\partial\mathcal L}{\partial\dot\phi(t,\mathbf{x})}
  =\dot\phi(t,\mathbf{x}).
\end{equation}
La densidad hamiltoniana es
\begin{equation}
  \eqbox{\mathcal{H}
  \coloneqq\frac{1}{2}\pi^2
  +\frac{1}{2}\abs{\grad\phi}^2
  +\frac{m^2}{2}\phi^2
  +\frac{\lambda}{4!}\phi^4.}
\end{equation}

Al integrar la densidad hamiltoniana sobre el espacio, separamos la parte libre
de la interacción:
\begin{equation}
  \boxed{\begin{aligned}
  H(t)
  &\coloneqq\int d^3x\,\mathcal{H}(t,\mathbf{x})
  \\
  &=
  \underbrace{\int d^3x\left[
    \frac{1}{2}\pi^2(t,\mathbf{x})
    +\frac{1}{2}\abs{\grad\phi(t,\mathbf{x})}^2
    +\frac{m^2}{2}\phi^2(t,\mathbf{x})
  \right]}_{\displaystyle H_0(t)\ \text{parte libre}}
  +
  \underbrace{\frac{\lambda}{4!}\int d^3x\,\phi^4(t,\mathbf{x})}
  _{\displaystyle H'(t)\ \text{interacción}} .
  \end{aligned}}
\end{equation}
Los argumentos temporales anteriores indican que se evalúan las variables
dinámicas en un instante; no representan una dependencia temporal explícita
del Hamiltoniano. Al pasar al formalismo de operadores, fijaremos la
separación $H=H_0+H'_S$ mediante los operadores iniciales en $t=0$.

\section{Imágenes dinámicas}

Usamos $\hbar=c=1$ y tomamos $t_0=0$. Consideramos una teoría cuyo Hamiltoniano
no depende explícitamente del tiempo.

\subsection{Campos y generadores de evolución}

\begin{definicion}[Campo y momento conjugado en la imagen de Schrödinger]
Fijamos los operadores iniciales del campo y de su momento conjugado:
\begin{equation}
  \eqbox{\phi_S(\mathbf{x})\coloneqq\phi(0,\mathbf{x}),
  \qquad
  \pi_S(\mathbf{x})\coloneqq\pi(0,\mathbf{x}).}
  \label{eq:phiSinitial}
\end{equation}
El subíndice $S$ indica la imagen de Schrödinger. Estos operadores no llevan
argumento temporal; representan los datos iniciales de la teoría completa.
Estas construcciones se usan en el formalismo canónico perturbativo;
los productos de campos en un mismo punto requieren regularización.
\end{definicion}


\begin{definicion}[Hamiltonianos en la imagen de Schrödinger]
Con estos operadores definimos la parte libre y el término de interacción:
\begin{equation}
  \eqbox{H_0\coloneqq\int d^3x\left[
    \frac{1}{2}\pi_S^2(\mathbf{x})
    +\frac{1}{2}\abs{\grad\phi_S(\mathbf{x})}^2
    +\frac{m^2}{2}\phi_S^2(\mathbf{x})\right].}
  \label{eq:H0Schrodinger}
\end{equation}
\begin{equation}
  \eqbox{H'_S\coloneqq\frac{\lambda}{4!}\int d^3x\,\phi_S^4(\mathbf{x}),
  \qquad H=H_0+H'_S.}
  \label{eq:Hprime}
\end{equation}
Así, $H$, $H_0$ y $H'_S$ son operadores fijos de Schrödinger.
\end{definicion}


\begin{definicion}[Operadores de evolución temporal]
Para los Hamiltonianos independientes del tiempo $H$ y $H_0$\footnote{\begin{minipage}[t]{\dimexpr\linewidth-1.8em\relax}
Para un Hamiltoniano fijo, la ecuación de Schrödinger tiene la solución
$\ket{\psi(t)}=e^{-iHt}\ket{\psi(0)}$. La condición $\partial H/\partial t=0$
significa que el Hamiltoniano no depende explícitamente del tiempo.
Esta idea ya aparece en ejemplos familiares de la mecánica clásica:
\begin{equation*}
\begin{aligned}
 H_{\mathrm{libre}}(\mathbf x,\mathbf p)
   &=\frac{\mathbf p^2}{2m}
   &&\text{(partícula libre)},\\
 H_{\mathrm{osc}}(x,p)
   &=\frac{p^2}{2m}+\frac12 kx^2
   &&\text{(oscilador armónico)},\\
 H_{\mathrm{Coul}}(\mathbf x,\mathbf p)
   &=\frac{\mathbf p^2}{2m}-\frac{\kappa}{r}
   &&\text{(potencial de Coulomb)},\\
 H_{\mathrm{grav}}(z,p_z)
   &=\frac{p_z^2}{2m}+mgz
   &&\text{(gravedad uniforme)},\\
 H_{\mathrm{pend}}(\theta,p_\theta)
   &=\frac{p_\theta^2}{2m\ell^2}+mg\ell(1-\cos\theta)
   &&\text{(péndulo simple)}.
\end{aligned}
\fnEqTag
\end{equation*}
Aquí $m$ es la masa, $k>0$ la constante del resorte,
$\kappa>0$ la intensidad de la atracción y $r=|\mathbf x|$.
En gravedad uniforme, $g$ es la aceleración gravitatoria, $z$ la altura
medida hacia arriba y $p_z$ su momento conjugado. Para el péndulo,
$\ell$ es la longitud fija, $\theta$ el ángulo respecto de la vertical hacia
abajo y $p_\theta=m\ell^2\dot\theta$ su momento conjugado; el punto de
suspensión permanece fijo. Suponemos que no hay rozamiento.
Con estos parámetros constantes, ninguno de los cinco Hamiltonianos contiene
$t$ explícitamente, aunque la posición y el momento cambien durante el
movimiento. En general, $H=\mathbf p^2/(2m)+V(\mathbf x)$ cumple esta
condición para cualquier potencial estático. Por tanto, un Hamiltoniano
independiente del tiempo puede describir movimiento e interacción.
\end{minipage}}, definimos los
operadores de evolución completa y libre, respectivamente, por
\begin{equation}
  U(t)\coloneqq e^{-iHt},
  \qquad
  U_0(t)\coloneqq e^{-iH_0t}.
\end{equation}
\end{definicion}

\begin{definicion}[Campo en la imagen de Heisenberg]
A partir del operador inicial $\phi_S(\mathbf{x})$, definimos el campo
interactuante en la imagen de Heisenberg mediante el Hamiltoniano completo
$H=H_0+H'_S$:
\begin{equation}
  \eqbox{\phi_H(t,\mathbf{x})\coloneqq
    e^{iHt}\phi_S(\mathbf{x})e^{-iHt}.}
  \label{eq:fieldHeisenberg}
\end{equation}
En la teoría $\lambda\phi^4$, este es el campo interactuante y describe la
evolución de la teoría completa.
\end{definicion}

\begin{definicion}[Campo en la imagen de interacción]
A partir del mismo operador inicial $\phi_S(\mathbf{x})$, definimos el campo
en la imagen de interacción mediante la parte libre $H_0$:
\begin{equation}
  \eqbox{\phi_I(t,\mathbf{x})\coloneqq
    e^{iH_0t}\phi_S(\mathbf{x})e^{-iH_0t}.}
  \label{eq:fieldInteraction}
\end{equation}
El subíndice $I$ indica la imagen de interacción; este campo tiene evolución
libre.\footnote{Seguimos la terminología de D. Tong,
\emph{Lectures on Quantum Field Theory}, secciones 3.1 y 3.7, donde se
distingue el campo en la imagen de Heisenberg de la teoría completa del campo
en la imagen de interacción. Véase
\href{https://www.damtp.cam.ac.uk/user/tong/qft/qfthtml/S3.html}{Interacting Fields}.}
\end{definicion}

\begin{observacion}
La igualdad \(\phi_S(\mathbf{x})=\phi_H(0,\mathbf{x})\) no supone que el campo
interactuante sea libre: \(\phi_H(0,\mathbf{x})\) es el campo completo de la
teoría \(\lambda\phi^4\) evaluado en el tiempo inicial. En \(t=0\), las
imágenes de Schrödinger, Heisenberg e interacción comparten el mismo operador de
campo,
\begin{equation}
  \phi_S(\mathbf{x})=\phi_H(0,\mathbf{x})=\phi_I(0,\mathbf{x}).
\end{equation}
Si ese dato inicial se evoluciona con \(H_0\), obtenemos el campo libre
\(\phi_I(t,\mathbf{x})\). Si se evoluciona con el Hamiltoniano completo
\(H=H_0+H'_S\), obtenemos el campo completo \(\phi_H(t,\mathbf{x})\). Por eso, en
general, \(\phi_H(t,\mathbf{x})\neq\phi_I(t,\mathbf{x})\) si \(t\neq0\).
\end{observacion}

\begin{center}
\begin{minipage}{\linewidth}
\centering
\begin{tikzpicture}[
  >=Latex,
  box/.style={draw,rounded corners=2pt,align=center,inner sep=6pt},
  evolution/.style={->,very thick,blue}
]
  \node[box] (initial) at (0,0)
    {$t=0$\\[2pt]
    $\phi_S(\mathbf{x})=\phi_H(0,\mathbf{x})=\phi_I(0,\mathbf{x})$\\
    dato inicial común};
  \node[box] (heis) at (7.0,1.35)
    {$\phi_H(t,\mathbf{x})=e^{iHt}\phi_S(\mathbf{x})e^{-iHt}$\\
    evolución completa};
  \node[box] (inter) at (7.0,-1.35)
    {$\phi_I(t,\mathbf{x})=e^{iH_0t}\phi_S(\mathbf{x})e^{-iH_0t}$\\
    evolución libre};
  \draw[evolution] (initial.east) --
    node[pos=0.68,above=3pt,sloped,fill=white,inner sep=1pt] {\scriptsize $H=H_0+H'_S$}
    (heis.west);
  \draw[evolution] (initial.east) --
    node[pos=0.58,below=3pt,sloped,fill=white,inner sep=1pt] {\scriptsize $H_0$}
    (inter.west);
\end{tikzpicture}
\diagramcaption{Resumen de \eqref{eq:fieldHeisenberg} y \eqref{eq:fieldInteraction}: el mismo
dato inicial evoluciona con $H$ para dar $\phi_H$ y con $H_0$ para dar $\phi_I$.}
\end{minipage}
\end{center}

\subsection{Estados y operadores en cada imagen}

\begin{definicion}[Imagen de Schrödinger]
Los estados evolucionan con el Hamiltoniano completo y los operadores sin dependencia
temporal explícita permanecen fijos:
\begin{equation}
  \eqbox{\ket{\psi_S(t)}\coloneqq e^{-iHt}\ket{\psi_S(0)},
  \qquad \frac{\partial A_S}{\partial t}=0.}
\end{equation}
Por tanto, el operador se denota $A_S$, sin argumento temporal; no se escribe $A_S(t)$.
\end{definicion}

\begin{definicion}[Imagen de Heisenberg]
Los estados son constantes y los operadores evolucionan con el Hamiltoniano completo:
\begin{equation}
  \eqbox{\ket{\psi_H}\coloneqq\ket{\psi_S(0)},
  \qquad
  A_H(t)\coloneqq e^{iHt}A_Se^{-iHt}.}
\end{equation}
\end{definicion}

\begin{definicion}[Imagen de interacción]
Los operadores evolucionan con $H_0$:
\begin{equation}
  \eqbox{A_I(t)
  \coloneqq e^{iH_0t}A_Se^{-iH_0t}.}
  \label{eq:AI}
\end{equation}

De la definición de la imagen de Heisenberg,
\begin{equation}
  A_H(t)=e^{iHt}A_Se^{-iHt},
\end{equation}
se despeja
\begin{equation}
  A_S=e^{-iHt}A_H(t)e^{iHt}.
\end{equation}
Al sustituir este resultado en \eqref{eq:AI},
\begin{equation}
  \eqbox{A_I(t)
  =e^{iH_0t}e^{-iHt}A_H(t)e^{iHt}e^{-iH_0t}.}
  \label{eq:AIfromAHexplicit}
\end{equation}
\begin{equation}
  \eqbox{A_I(t)=U_I(t)A_H(t)U_I^\dagger(t).}
  \label{eq:AIfromAH}
\end{equation}
donde
\begin{equation}
  \eqbox{U_I(t)\coloneqq e^{iH_0t}e^{-iHt}.}
  \label{eq:UI}
\end{equation}
Aplicada al campo, la relación \eqref{eq:AIfromAHexplicit} muestra que la
transformación a la imagen de interacción cancela la evolución completa y deja
sólo la evolución libre:
\begin{align}
  \phi_I(t,\mathbf{x})
  &=e^{iH_0t}e^{-iHt}\phi_H(t,\mathbf{x})e^{iHt}e^{-iH_0t}
  \notag\\
  &=e^{iH_0t}e^{-iHt}
    \left(e^{iHt}\phi_S(\mathbf{x})e^{-iHt}\right)
    e^{iHt}e^{-iH_0t}
  \notag\\
  &=e^{iH_0t}\phi_S(\mathbf{x})e^{-iH_0t}.
  \label{eq:phiIfromphiH}
\end{align}
Equivalentemente,
\begin{equation}
  \eqbox{A_H(t)=U_I^\dagger(t)A_I(t)U_I(t).}
\end{equation}

Los estados contienen la evolución restante:
\begin{equation}
  \ket{\psi_I(t)}\coloneqq e^{iH_0t}\ket{\psi_S(t)},
  \qquad
  i\frac{d}{dt}\ket{\psi_I(t)}=H'_I(t)\ket{\psi_I(t)},
\end{equation}
\end{definicion}



\begin{center}
\begin{minipage}{\linewidth}
\centering
\begin{tikzpicture}
\begin{axis}[
 width=12.5cm,height=8.9cm,view={-35}{27},
 axis lines=none,ticks=none,
 xmin=-3,xmax=3,ymin=0,ymax=4,zmin=0,zmax=1.35,
 clip=false,z buffer=sort,
 colormap={campo}{color(0cm)=(white);color(1cm)=(diagramblue!38!white)}]
% x es posición, y es tiempo y z es valor de un perfil ilustrativo.
\addplot3[surf,shader=faceted interp,faceted color=diagramblue!28,
 domain=-3:3,y domain=0:4,samples=35,samples y=23]
 {exp(-1.25*(x-.4*y+.8)^2)};
% El perfil inicial queda destacado y es el borde de la superficie.
\addplot3[diagramgreen,very thick,domain=-3:3,samples=100,samples y=1]
 (x,0,{exp(-1.25*(x+.8)^2)});
% Dos perfiles posteriores muestran cortes espaciales de la evolución.
\addplot3[diagramblue,thick,domain=-3:3,samples=100,samples y=1]
 (x,2,{exp(-1.25*x^2)});
\addplot3[diagramblue,thick,domain=-3:3,samples=100,samples y=1]
 (x,4,{exp(-1.25*(x-.8)^2)});
% Los tres ejes tienen significados distintos.
\draw[gray!65,->,>=stealth,thin] (axis cs:-3.25,0,0)--(axis cs:3.6,0,0)
 node[below right,text=black] {$x$};
\draw[diagramred,->,>=stealth,thick] (axis cs:3.3,0,0)--(axis cs:3.3,4.7,0)
 node[above,text=diagramred] {$t$};
\draw[gray!65,->,>=stealth,thin] (axis cs:-3.25,0,-.18)--(axis cs:-3.25,0,2.25)
 node[above,text=black] {$\varphi$};
\node[anchor=north,text=diagramgreen,font=\small,align=center]
 at (axis cs:-.8,0,-.16) {perfil inicial\\$\varphi(0,x)$};
\node[anchor=west,text=diagramblue,font=\small]
 at (axis cs:1.3,4,1.18) {$\varphi(t,x)$};
\draw[diagramblue!65,thin] (axis cs:1.3,4,1.15)--(axis cs:.8,4,1);
\node[text=diagramred,font=\small,rotate=27]
 at (axis cs:3.65,2,.06) {evolución};
\node[text=gray!75,font=\footnotesize]
 at (axis cs:-2.8,4,.08) {$t>0$};
\end{axis}
\end{tikzpicture}
\diagramcaption{Del perfil espacial a su evolución temporal. La curva verde
representa un perfil inicial en $t=0$; los cortes azules muestran perfiles
posteriores. El eje rojo indica el tiempo y la altura indica el valor del
campo. La superficie es un esquema de un perfil clásico, no una gráfica
del operador cuántico ni una solución calculada de la teoría $\lambda\phi^4$.}
\end{minipage}
\end{center}

La figura ilustra cómo un perfil espacial se extiende a lo largo del tiempo.
En las definiciones operatoriales, ese mismo tiempo aparece tanto en $\phi_I$
como en $\phi_H$; sus evoluciones se distinguen por el generador, $H_0$ o $H$.

Un ejemplo concreto en una dimensión espacial es el modo libre
\[
 \varphi(t,x)=A\cos(kx-\omega_k t),\qquad
 \omega_k=\sqrt{k^2+m^2}.
\]
En $t=0$ tenemos el perfil $A\cos(kx)$; a tiempos posteriores,
la fase $-\omega_k t$ desplaza las crestas y los valles. Esta onda es
un ejemplo analítico de evolución libre; la superficie del dibujo
sólo ilustra la construcción geométrica de un perfil a lo largo del tiempo.

\section{Hamiltoniano de interacción}

Para expresar $H'_I(t)$ en términos de campos libres, comprobamos primero
mediante la expansión modal la evolución con $H_0$ definida en
\eqref{eq:fieldInteraction}.

\subsection{Reconstrucción del campo libre}

\begin{ejemplo}[Comprobación de la evolución libre]
Escribimos el operador inicial común $\phi_S(\mathbf{x})=\phi_I(0,\mathbf{x})$
en la expansión modal asociada al Hamiltoniano libre $H_0$:
\begin{equation}
  \phi_S(\mathbf{x})=
  \int\frac{d^3k}{(2\pi)^3}\frac{1}{\sqrt{2\omega_{\mathbf{k}}}}
  \left[
    a_{\mathbf{k}}e^{i\mathbf{k}\cdot\mathbf{x}}
    +a_{\mathbf{k}}^\dagger e^{-i\mathbf{k}\cdot\mathbf{x}}
  \right].
  \label{eq:freefield0}
\end{equation}
Aquí $\omega_{\mathbf{k}}\coloneqq\sqrt{\mathbf{k}^2+m^2}$.
Los operadores $a_{\mathbf{k}}$ y $a_{\mathbf{k}}^\dagger$ se definen con
respecto a esta parte libre, cuyo vacío satisface $a_{\mathbf{k}}\ket{0}=0$.
Considere el Hamiltoniano libre, omitiendo la constante de energía de punto cero,
\begin{equation}
  H_0=\int\frac{d^3k}{(2\pi)^3}\,
  \omega_{\mathbf{k}}a_{\mathbf{k}}^\dagger a_{\mathbf{k}},
\end{equation}
y las relaciones de conmutación
\begin{equation}
  [a_{\mathbf{k}},a_{\mathbf{p}}^\dagger]
  =(2\pi)^3\delta^{(3)}(\mathbf{k}-\mathbf{p}),
  \qquad
  [a_{\mathbf{k}},a_{\mathbf{p}}]
  =[a_{\mathbf{k}}^\dagger,a_{\mathbf{p}}^\dagger]=0.
\end{equation}
Demuestre, mediante los conmutadores anidados y la fórmula de
Baker--Campbell--Hausdorff, que
\begin{equation}
  \boxed{\begin{aligned}
  e^{iH_0t}\phi_S(\mathbf{x})e^{-iH_0t}
  &={}\int\frac{d^3k}{(2\pi)^3}\frac{1}{\sqrt{2\omega_{\mathbf{k}}}}
  \left[
    a_{\mathbf{k}}e^{-i\omega_{\mathbf{k}}t+i\mathbf{k}\cdot\mathbf{x}}
    +a_{\mathbf{k}}^\dagger e^{i\omega_{\mathbf{k}}t-i\mathbf{k}\cdot\mathbf{x}}
  \right].
  \end{aligned}}
  \label{eq:freefield}
\end{equation}
Concluya que la evolución generada por $H_0$ reconstruye el campo libre de
Klein--Gordon $\phi_{\mathrm{libre}}(t,\mathbf{x})$.

\par\smallskip\noindent\textbf{Solución.}
Según la definición de la imagen de interacción, calculamos
\begin{equation}
  \eqbox{\phi_I(t,\mathbf{x})
  =e^{iH_0t}\phi_S(\mathbf{x})e^{-iH_0t}.}
  \label{eq:fieldBCH}
\end{equation}
Primero obtenemos los conmutadores elementales\footnote{Usamos
\([AB,C]=A[B,C]+[A,C]B\). Entonces
\begin{align*}
[H_0,a_{\mathbf{k}}]
&=\int\frac{d^3p}{(2\pi)^3}\,
\omega_{\mathbf{p}}
[a_{\mathbf{p}}^\dagger a_{\mathbf{p}},a_{\mathbf{k}}] \\
&=\int\frac{d^3p}{(2\pi)^3}\,
\omega_{\mathbf{p}}
\bigl(a_{\mathbf{p}}^\dagger[a_{\mathbf{p}},a_{\mathbf{k}}]
+[a_{\mathbf{p}}^\dagger,a_{\mathbf{k}}]a_{\mathbf{p}}\bigr) \\
&=-\int\frac{d^3p}{(2\pi)^3}\,
\omega_{\mathbf{p}}(2\pi)^3
\delta^{(3)}(\mathbf{p}-\mathbf{k})a_{\mathbf{p}}
=-\omega_{\mathbf{k}}a_{\mathbf{k}}.
\fnEqTag
\end{align*}
y de manera análoga
\begin{align*}
[H_0,a_{\mathbf{k}}^\dagger]
&=\int\frac{d^3p}{(2\pi)^3}\,
\omega_{\mathbf{p}}
[a_{\mathbf{p}}^\dagger a_{\mathbf{p}},a_{\mathbf{k}}^\dagger] \\
&=\int\frac{d^3p}{(2\pi)^3}\,
\omega_{\mathbf{p}}
a_{\mathbf{p}}^\dagger[a_{\mathbf{p}},a_{\mathbf{k}}^\dagger]
=\omega_{\mathbf{k}}a_{\mathbf{k}}^\dagger .
\fnEqTag
\end{align*}
}
\begin{equation}
  [H_0,a_{\mathbf{k}}]=-\omega_{\mathbf{k}}a_{\mathbf{k}},
  \qquad
  [H_0,a_{\mathbf{k}}^\dagger]=\omega_{\mathbf{k}}a_{\mathbf{k}}^\dagger.
  \label{eq:H0akcomm}
\end{equation}

Ahora aplicamos la fórmula de conjugación de Baker--Campbell--Hausdorff
(también llamada lema de Hadamard) con
$A=iH_0t$ y $B=\phi_S(\mathbf{x})=\phi_I(0,\mathbf{x})$:
\begin{equation}
  \boxed{\begin{gathered}\textstyle
  e^{iH_0t}B e^{-iH_0t}
  =
  B+it[H_0,B]
  +\frac{(it)^2}{2!}[H_0,[H_0,B]]
  +\frac{(it)^3}{3!}[H_0,[H_0,[H_0,B]]]+\cdots .
  \end{gathered}}
  \label{eq:BCHfield}
\end{equation}
Calculemos entonces los conmutadores que aparecen en \eqref{eq:BCHfield}.
Usando \eqref{eq:H0akcomm} en la expansión \eqref{eq:freefield0}, el primer
conmutador es
\begin{equation}
  [H_0,\phi_S(\mathbf{x})]
  =
  \int\frac{d^3k}{(2\pi)^3}
  \frac{1}{\sqrt{2\omega_{\mathbf{k}}}}
  \left[
    -\omega_{\mathbf{k}}a_{\mathbf{k}}e^{i\mathbf{k}\cdot\mathbf{x}}
    +\omega_{\mathbf{k}}a_{\mathbf{k}}^\dagger e^{-i\mathbf{k}\cdot\mathbf{x}}
  \right].
  \label{eq:firstPhiComm}
\end{equation}
Al conmutar una vez más con \(H_0\),
\begin{equation}
  [H_0,[H_0,\phi_S(\mathbf{x})]]
  =
  \int\frac{d^3k}{(2\pi)^3}
  \frac{1}{\sqrt{2\omega_{\mathbf{k}}}}
  \left[
    \omega_{\mathbf{k}}^2a_{\mathbf{k}}e^{i\mathbf{k}\cdot\mathbf{x}}
    +\omega_{\mathbf{k}}^2a_{\mathbf{k}}^\dagger e^{-i\mathbf{k}\cdot\mathbf{x}}
  \right].
  \label{eq:secondPhiComm}
\end{equation}
El tercer conmutador ya deja ver la alternancia de signos:
\begin{equation}
  [H_0,[H_0,[H_0,\phi_S(\mathbf{x})]]]
  =
  \int\frac{d^3k}{(2\pi)^3}
  \frac{1}{\sqrt{2\omega_{\mathbf{k}}}}
  \left[
    -\omega_{\mathbf{k}}^3a_{\mathbf{k}}e^{i\mathbf{k}\cdot\mathbf{x}}
    +\omega_{\mathbf{k}}^3a_{\mathbf{k}}^\dagger e^{-i\mathbf{k}\cdot\mathbf{x}}
  \right].
  \label{eq:thirdPhiComm}
\end{equation}
Por tanto, para \(n\) conmutadores anidados,
\begin{equation}
  \underbrace{[H_0,[H_0,\cdots[H_0,\phi_S(\mathbf{x})]\cdots]]}_{n\ \text{veces}}
  =
  \int\frac{d^3k}{(2\pi)^3}
  \frac{1}{\sqrt{2\omega_{\mathbf{k}}}}
  \left[
    (-\omega_{\mathbf{k}})^n
    a_{\mathbf{k}}e^{i\mathbf{k}\cdot\mathbf{x}}
    +\omega_{\mathbf{k}}^n
    a_{\mathbf{k}}^\dagger e^{-i\mathbf{k}\cdot\mathbf{x}}
  \right].
  \label{eq:nestedPhiComm}
\end{equation}
Sustituyendo estos conmutadores en \eqref{eq:BCHfield},
\begin{align}
  \phi_I(t,\mathbf{x})
  &\coloneqq e^{iH_0t}\phi_S(\mathbf{x})e^{-iH_0t}
  \notag\\
  &=
  \int\frac{d^3k}{(2\pi)^3}
  \frac{1}{\sqrt{2\omega_{\mathbf{k}}}}
  \left[
    a_{\mathbf{k}}e^{i\mathbf{k}\cdot\mathbf{x}}
    \sum_{n=0}^{\infty}\frac{(-i\omega_{\mathbf{k}}t)^n}{n!}
    +
    a_{\mathbf{k}}^\dagger e^{-i\mathbf{k}\cdot\mathbf{x}}
    \sum_{n=0}^{\infty}\frac{(i\omega_{\mathbf{k}}t)^n}{n!}
  \right]
  \notag\\
  &=
  \int\frac{d^3k}{(2\pi)^3}
  \frac{1}{\sqrt{2\omega_{\mathbf{k}}}}
  \left[
    a_{\mathbf{k}}e^{-i\omega_{\mathbf{k}}t+i\mathbf{k}\cdot\mathbf{x}}
    +
    a_{\mathbf{k}}^\dagger
    e^{i\omega_{\mathbf{k}}t-i\mathbf{k}\cdot\mathbf{x}}
  \right]
  \equiv\phi_{\mathrm{libre}}(t,\mathbf{x}).
  \label{eq:BCHsubstitutionField}
\end{align}
La última igualdad define la notación $\phi_{\mathrm{libre}}(t,\mathbf{x})$
para el campo libre de Klein--Gordon dado por esta expansión modal. Por tanto,
\begin{equation}
  \eqbox{\phi_I(t,\mathbf{x})=\phi_{\mathrm{libre}}(t,\mathbf{x}).}
  \label{eq:BCHrecoversField}
\end{equation}
La suma de los conmutadores reconstruye las fases $e^{\mp i\omega_{\mathbf{k}}t}$
y verifica la expansión modal del campo libre anunciada en el enunciado.
\end{ejemplo}

\begin{observacion}
El campo $\phi_I$ no es una ``parte'' de $\phi_H$ en sentido aditivo.
Es el campo libre con el que organizamos perturbativamente el efecto de
$H'_I(t)$.
\end{observacion}

\subsection{Construcción de la interacción}

Una vez calculado el campo $\phi_I$, transformamos un operador distinto:
el término de interacción $H'_S$ definido en \eqref{eq:Hprime}. Recordemos que
\begin{equation}
  H'_S=\frac{\lambda}{4!}\int d^3z\,\phi_S^4(\mathbf{z}).
\end{equation}
Su expresión en la imagen de interacción es
\begin{equation}
  \eqbox{H'_I(t)\coloneqq e^{iH_0t}H'_Se^{-iH_0t}.}
  \label{eq:HIdefinition}
\end{equation}
Usamos ahora el resultado del Ejemplo 1, \eqref{eq:BCHrecoversField}:
$e^{iH_0t}\phi_S(\mathbf{z})e^{-iH_0t}=\phi_I(t,\mathbf{z})$.
La conjugación conserva los productos: insertamos
$e^{-iH_0t}e^{iH_0t}=1$ entre los cuatro factores de $\phi_S^4$ y aplicamos
este resultado a cada factor. Sustituyendo \eqref{eq:Hprime},
\begin{align}
  e^{iH_0t}H'_Se^{-iH_0t}
  &=\frac{\lambda}{4!}\int d^3z\,
    e^{iH_0t}\phi_S^4(\mathbf{z})e^{-iH_0t}
    \notag\\
  &=\frac{\lambda}{4!}\int d^3z\,
    \left(e^{iH_0t}\phi_S(\mathbf{z})e^{-iH_0t}\right)^4
    \notag\\
  &=\frac{\lambda}{4!}\int d^3z\,\phi_I^4(t,\mathbf{z}).
  \label{eq:HSaHI}
\end{align}
Por tanto,
\begin{equation}
  \eqbox{H'_I(t)
  =e^{iH_0t}H'_Se^{-iH_0t}
  =\frac{\lambda}{4!}\int d^3z\,\phi_I^4(t,\mathbf{z}).}
  \label{eq:HIphi4}
\end{equation}

Aquí aparece sólo el término de interacción porque partimos de $H'_S$.
Al transformar el Hamiltoniano total, la parte libre también está presente:
\begin{equation}
  e^{iH_0t}He^{-iH_0t}=H_0+H'_I(t).
  \label{eq:totalHamiltonianI}
\end{equation}
Para los estados, en cambio, el cambio de imagen absorbe la evolución libre.
Al derivar $\ket{\psi_I(t)}=e^{iH_0t}\ket{\psi_S(t)}$, obtenemos
\begin{align}
  i\frac{d}{dt}\ket{\psi_I(t)}
  &=\left[-H_0+e^{iH_0t}He^{-iH_0t}\right]\ket{\psi_I(t)}
  \notag\\
  &=H'_I(t)\ket{\psi_I(t)}.
  \label{eq:interactionGenerator}
\end{align}
El término $-H_0$ procede de derivar $e^{iH_0t}$ y cancela la parte libre de
\eqref{eq:totalHamiltonianI}. Por eso la serie de Dyson, que describe la
evolución de estos estados, utiliza $H'_I(t)$.

La interacción se usa sin orden normal: el término $\phi_I^4$ no se
reemplaza por $:\!\phi_I^4\!:$, de modo que se conservan las contracciones
locales que originan los lazos de los ejemplos.

La integral de \eqref{eq:HIphi4} recorre todo el espacio a tiempo fijo. Por eso
$H'_I(t)$ sigue dependiendo de $t$. El tiempo se integra posteriormente en la serie de
Dyson:
\begin{equation}
  \eqbox{\int_{-\infty}^{\infty}dt\,H'_I(t)
  =\frac{\lambda}{4!}\int d^4z\,\phi_I^4(z),
  \qquad z\coloneqq(t,\mathbf{z}).}
  \label{eq:intHI}
\end{equation}

\section{Punto de partida del cálculo perturbativo}

Una vez construidos $\phi_I$ y $H'_I$, nuestros cálculos comienzan con la
fórmula de Gell-Mann--Low\footnote{Murray Gell-Mann y Francis E. Low
introdujeron este resultado en el contexto de la formulación perturbativa de
teorías cuánticas de campos. En estas notas lo usamos como la regla que permite
escribir valores esperados en el vacío interactuante $\ket{\Omega}$ como
cocientes de valores esperados en el vacío libre $\ket{0}$.} para el correlador
de dos puntos.
De hecho, esta fórmula funciona precisamente porque reemplaza correladores del
campo completo por correladores del campo libre de Klein--Gordon, junto con la
inserción ordenada temporalmente de la interacción.
\begin{equation}
  \boxed{\displaystyle
  \mel{\Omega}{T\{\phi(x)\phi(y)\}}{\Omega}
  =
  \frac{
    \mel{0}{T\{\phi_I(x)\phi_I(y)
    e^{-i\int_{-\infty}^{\infty}dt\,H'_I(t)}\}}{0}
  }{
    \mel{0}{T\{e^{-i\int_{-\infty}^{\infty}dt\,H'_I(t)}\}}{0}
  }
  \coloneqq\frac{N}{D}.}
  \label{eq:GML}
\end{equation}
Aquí $\phi(x)$ denota $\phi_H(x)$, el campo completo en la imagen de
Heisenberg; $\ket{\Omega}$ es el vacío interactuante y $\ket{0}$ el vacío
libre, ambos normalizados. Los símbolos $N$ y $D$ denotan respectivamente
el numerador y el denominador de \eqref{eq:GML}.
La fórmula se utiliza perturbativamente, con la preparación del vacío y
los límites temporales definidos mediante una prescripción de convergencia
compatible con las condiciones de Feynman. Se supone que el vacío libre
tiene solapamiento no nulo con el vacío interactuante en la teoría
regularizada. El $i\epsilon$ fija la prescripción de los polos de los
propagadores; no elimina las divergencias ultravioletas, que requieren
regularización y renormalización.

Sustituyendo \eqref{eq:HIphi4} en \eqref{eq:GML}, el operador de evolución
contiene, por \eqref{eq:intHI}, la exponencial de
$-i\frac{\lambda}{4!}\int d^4z\,\phi_I^4(z)$ bajo el ordenamiento temporal.
Primero expandimos formalmente la exponencial en su serie de Taylor:
\begin{equation}
  \begin{aligned}
  \exp\left[-i\frac{\lambda}{4!}\int d^4z\,\phi_I^4(z)\right]
  &=1-i\frac{\lambda}{4!}\int d^4z\,\phi_I^4(z)
  \\
  &\quad+\frac{1}{2!}\left(-i\frac{\lambda}{4!}\right)^2
  \left[\int d^4z\,\phi_I^4(z)\right]^2+\cdots .
  \end{aligned}
  \label{eq:DysonTaylor}
\end{equation}
En el término cuadrático, cada factor de la integral tiene su propia variable
de integración. Como $z$ es una variable muda, podemos llamarlas $z_1$ y $z_2$:
\begin{equation}
  \begin{aligned}
  \left[\int d^4z\,\phi_I^4(z)\right]^2
  &=\left[\int d^4z_1\,\phi_I^4(z_1)\right]
    \left[\int d^4z_2\,\phi_I^4(z_2)\right]
  \\
  &=\int d^4z_1\int d^4z_2\,
    \phi_I^4(z_1)\phi_I^4(z_2).
  \end{aligned}
  \label{eq:DysonIndependentVariables}
\end{equation}
Los puntos $z_1$ y $z_2$ se integran independientemente; el cuadrado de la
integral produce campos evaluados en ambos puntos, no ocho campos en un único
punto. A orden $n$ aparecen $n$ variables $z_1,\ldots,z_n$, una por cada
factor de la integral. Aplicamos ahora $T$ a los campos de cada término,
ordenándolos según sus tiempos $z_j=(t_j,\mathbf{z}_j)$, y obtenemos
\begin{equation}
  \boxed{\begin{aligned}
  T\left\{\exp\left[-i\frac{\lambda}{4!}
  \int d^4z\,\phi_I^4(z)\right]\right\}
  &=1-i\frac{\lambda}{4!}\int d^4z_1\,
  T\{\phi_I^4(z_1)\}
  \\
  &\quad
  +\frac{(-i\lambda)^2}{2!(4!)^2}
  \int d^4z_1\int d^4z_2\,
  T\{\phi_I^4(z_1)\phi_I^4(z_2)\}+\cdots .
  \end{aligned}}
  \label{eq:Dyson}
\end{equation}

Por tanto,
\begin{align}
  N
  &=\mel{0}{T\{\phi_I(x)\phi_I(y)\}}{0}
  -i\frac{\lambda}{4!}\int d^4z\,
  \mel{0}{T\{\phi_I(x)\phi_I(y)\phi_I^4(z)\}}{0}
  +O(\lambda^2),\\
  D
  &=1-i\frac{\lambda}{4!}\int d^4z\,
  \mel{0}{T\{\phi_I^4(z)\}}{0}
  +O(\lambda^2).
\end{align}
El denominador $D$ contiene las burbujas de vacío. Al formar el cociente
$N/D$, esas piezas desconectadas se cancelan, de modo que finalmente quedan
las contribuciones conectadas al par externo $x,y$.

\section{Teorema de Wick a primer orden}

La contracción básica es el propagador libre
\begin{equation}
  \Delta_F(a-b)
  \coloneqq\mel{0}{T\{\phi_I(a)\phi_I(b)\}}{0}.
\end{equation}
Las contracciones del término con seis campos se agrupan como
\begin{equation}
  \boxed{\begin{aligned}
  \mel{0}{T\{\phi_I(x)\phi_I(y)\phi_I^4(z)\}}{0}
  &=3\,\Delta_F(x-y)\bigl[\Delta_F(z-z)\bigr]^2
  \\
  &\quad
  +12\,\Delta_F(x-z)\Delta_F(y-z)\Delta_F(z-z).
  \end{aligned}}
  \label{eq:Wick6}
\end{equation}
El factor $3$ corresponde a las contracciones en las que $x$ se une con $y$; el factor
$12$ corresponde a las contracciones en las que ambos puntos externos se unen al vértice.

Para el denominador, Wick permite emparejar los cuatro campos en $z$ de tres
maneras. Si etiquetamos provisionalmente los campos como $1,2,3,4$, los
emparejamientos son $(12)(34)$, $(13)(24)$ y $(14)(23)$. Todos dan el mismo
producto de dos propagadores:
\begin{equation}
  \mel{0}{T\{\phi_I^4(z)\}}{0}
  =3\bigl[\Delta_F(z-z)\bigr]^2,
  \qquad \frac{3}{4!}=\frac{1}{8}.
  \label{eq:WickVacuumFour}
\end{equation}
De \eqref{eq:Wick6} y \eqref{eq:WickVacuumFour},
\begin{align}
  N
  &=\Delta_F(x-y)
  -i\lambda\int d^4z\left[
    \frac{1}{8}\Delta_F(x-y)\bigl[\Delta_F(z-z)\bigr]^2
  \right.
  \notag\\
  &\hspace{5.4cm}\left.
    +\frac{1}{2}\Delta_F(x-z)\Delta_F(y-z)\Delta_F(z-z)
  \right]+O(\lambda^2),\\
  D
  &=1-i\frac{\lambda}{8}\int d^4z\,
  \bigl[\Delta_F(z-z)\bigr]^2+O(\lambda^2).
\end{align}

\begin{center}
\begin{minipage}{\linewidth}
\centering
  \begin{tikzpicture}[scale=0.65,baseline=(current bounding box.center)]
    \draw[blue,very thick] (0,0)
      .. controls (-1.1,1.8) and (1.1,1.8) .. (0,0);
    \draw[blue,very thick] (0,0)
      .. controls (-1.1,-1.8) and (1.1,-1.8) .. (0,0);
    \fill[red] (0,0) circle (2.2pt);
    \node[right=5pt,text=black] at (0,0) {$z$};
    \draw[->,thick] (-0.48,1.75)
      .. controls (-1.05,2.25) and (1.05,2.25) .. (0.48,1.75);
    \node[above] at (0,2.1) {giro del lazo: $1/2$};
    \draw[->,thick] (-0.48,-1.75)
      .. controls (-1.05,-2.25) and (1.05,-2.25) .. (0.48,-1.75);
    \node[below] at (0,-2.1) {giro del lazo: $1/2$};
    \draw[<->,thick] (1.25,1.1)
      .. controls (2.15,0.8) and (2.15,-0.8) .. (1.25,-1.1);
    \node[right,align=left] at (2.05,0)
      {intercambio\\de lazos: $1/2$};
  \end{tikzpicture}
  \diagramcaption{Burbuja de vacío en figura de ocho: dos lazos unidos en un vértice.}
\end{minipage}
\end{center}
Su factor de simetría es $S=2\cdot2\cdot2=8$\footnote{Véase K.~Peeters y
M.~Zamaklar, \textit{A guide to quantum field theory}, sección 4.3,
discusión de la ecuación (4.34), pp.~31--32: se cuentan las inversiones de cada lazo y el intercambio de
los dos lazos de la figura de ocho.
\href{https://maths.dur.ac.uk/users/kasper.peeters/pdf/qft2.pdf}{Notas de Durham}.}:
podemos intercambiar los dos
extremos de cada lazo y también intercambiar los lazos entre sí, sin cambiar
el diagrama. Cada una de estas simetrías de orden dos aporta un factor $1/2$:
\begin{equation}
  \frac{1}{S}=\frac{1}{2}\,\frac{1}{2}\,\frac{1}{2}=\frac{1}{8},
  \qquad
  \mathcal{A}_{\mathrm{vac}}^{(1)}
  =\frac{-i\lambda}{8}\int d^4z\,\bigl[\Delta_F(z-z)\bigr]^2.
\end{equation}
El intercambio de los extremos de un lazo suele visualizarse como un giro
del lazo; lo que se cuenta es esa permutación discreta, no cualquier rotación
del dibujo. Se trata de una sola burbuja de vacío con dos lazos, no de dos
burbujas independientes. El $1/8$ obtenido por Wick ya incluye estas
simetrías: no se divide de nuevo por $8$.

\subsection{Cancelación de las burbujas de vacío}
La cancelación no ocurre únicamente a primer orden. Cada diagrama del
numerador se separa en una parte ligada a los puntos externos $x,y$ y un
conjunto de burbujas de vacío desconectadas de ellos. Estas burbujas forman
exactamente el mismo factor que aparece en el denominador. Diagramáticamente,
la normalización se escribe

\begin{equation}
  \boxed{\begin{aligned}
  \frac{N}{D}
  &=\frac{
    \left[\qftFreeDiagram+\qftTadpoleDiagram+\qftChainDiagram+\cdots\right]
    \exp\left(\qftVacuumDiagram+\cdots\right)
  }{\exp\left(\qftVacuumDiagram+\cdots\right)}
  \\[0.8em]
  &=\qftFreeDiagram+\qftTadpoleDiagram+\qftChainDiagram+\cdots .
  \end{aligned}}
  \label{eq:VacuumDiagramCancellation}
\end{equation}
\begin{center}
\begin{minipage}{\linewidth}
\centering
  \diagramcaption{Cancelación de las burbujas de vacío en el correlador normalizado.}
\end{minipage}
\end{center}

Cada dibujo representa su contribución completa, con los propagadores,
vértices, integrales y factores de simetría correspondientes. Los extremos
de las líneas abiertas son $x,y$. La exponencial contiene la suma de
burbujas de vacío conectadas; al expandirla, genera también conjuntos de
burbujas independientes, con sus factores $1/n!$.

Tras cancelar el factor común de burbujas de vacío en
\eqref{eq:VacuumDiagramCancellation}, queda la suma de diagramas conectados
al par externo $x,y$. Suponemos un vacío que conserva la simetría
$\phi\mapsto-\phi$, por lo que los correladores de un solo campo se anulan.
Los lazos unidos a los puntos externos, como el del \emph{tadpole},
permanecen en esta suma.

Para obtener el correlador hasta primer orden, conservamos únicamente
los dos primeros diagramas: la propagación libre, de orden $\lambda^0$,
y el \emph{tadpole}, de orden $\lambda$. La cadena y los demás términos
de segundo orden o superior quedan incluidos en $O(\lambda^2)$:
\begin{equation}
  \eqbox{\frac{N}{D}
  =\qftFreeDiagram+\qftTadpoleDiagram+O(\lambda^2).}
  \label{eq:VacuumFactorCancellation}
\end{equation}
Escribiendo las contribuciones de estos dos diagramas en espacio de
posiciones, obtenemos directamente
\begin{equation}
  \eqbox{\mel{\Omega}{T\{\phi(x)\phi(y)\}}{\Omega}
  =\Delta_F(x-y)
  -i\frac{\lambda}{2}\int d^4z\,
  \Delta_F(x-z)\Delta_F(y-z)\Delta_F(z-z)
  +O(\lambda^2).}
  \label{eq:G2primerorden}
\end{equation}
El término de orden cero tiene $S=1$ y el \emph{tadpole} de primer orden,
$S=2$. Las tres topologías de segundo orden se desarrollan en la sección
siguiente.
Aquí $\Delta_F(z-z)=\Delta_F(0)$ es formalmente divergente en el ultravioleta y debe
regularizarse y renormalizarse.\footnote{\begin{minipage}[t]{\dimexpr\linewidth-1.8em\relax}
Con nuestra convención,
\begin{equation*}
 \Delta_F(0)=\int\frac{d^3p}{(2\pi)^3}
 \frac{1}{2\sqrt{\mathbf p^2+m^2}}.
 \fnEqTag
\end{equation*}
Para momentos grandes, la medida radial convierte la integral en
una expresión que crece como $\int p\,dp$: de ahí la divergencia ultravioleta,
asociada a longitudes de onda pequeñas.
Un corte $|\mathbf p|\leq\Lambda$ regulariza la integral:
\begin{equation*}
 \Delta_F^{(\Lambda)}(0)=\frac{1}{4\pi^2}
 \int_0^\Lambda\frac{p^2\,dp}{\sqrt{p^2+m^2}}
 \sim\frac{\Lambda^2}{8\pi^2},\qquad \Lambda\gg m.
 \fnEqTag
\end{equation*}
La renormalización ajusta parámetros y contraterminos mediante condiciones
físicas para obtener predicciones finitas al retirar el regulador.
En este tadpole, la divergencia se absorbe en el contratermino de masa;
no es una burbuja de vacío que se cancele ni cambia $S=2$.
Véanse Peeters y Zamaklar,
\href{https://maths.dur.ac.uk/users/kasper.peeters/pdf/qft2.pdf}
{\textit{A guide to quantum field theory}}, secciones 7.1 y 7.2.\end{minipage}}
Esta divergencia no modifica los factores combinatorios
que se están obteniendo.

\section{Contribuciones conectadas de segundo orden}

Para la interacción $\lambda\phi^4/4!$ sin orden normal, las tres topologías
conectadas del correlador de dos puntos a segundo orden son las siguientes:

Dos diagramas \emph{tadpole} (renacuajos) en cadena:

\begin{center}
\begin{minipage}{\linewidth}
\centering
  \begin{tikzpicture}[scale=0.9,baseline=(current bounding box.center)]
    \draw[blue,very thick] (-2.5,0)--(2.5,0);
    \draw[blue,very thick] (-0.85,0)
      .. controls (-1.25,0.85) and (-0.45,0.85) .. (-0.85,0);
    \draw[blue,very thick] (0.85,0)
      .. controls (0.45,0.85) and (1.25,0.85) .. (0.85,0);
    \fill[red] (-0.85,0) circle (2.2pt) node[below=3pt,text=black] {$z_1$};
    \fill[red] (0.85,0) circle (2.2pt) node[below=3pt,text=black] {$z_2$};
    \node[below=3pt] at (-2.5,0) {$x$};
    \node[below=3pt] at (2.5,0) {$y$};
  \end{tikzpicture}
  \diagramcaption{Dos diagramas \emph{tadpole} en cadena.}
\end{minipage}
\end{center}
\begin{align}
  &\frac{(-i\lambda)^2}{4}
  \int d^4z_1\int d^4z_2\,
  \Delta_F(x-z_1)\Delta_F(z_1-z_1)
  \notag\\[-0.2em]
  &\hspace{3.2cm}\times
  \Delta_F(z_1-z_2)\Delta_F(z_2-z_2)\Delta_F(z_2-y).
  \label{eq:G2chainContribution}
\end{align}


Dos líneas entre vértices y un lazo local:
\begin{center}
\begin{minipage}{\linewidth}
\centering
  \begin{tikzpicture}[scale=0.9,baseline=(current bounding box.center)]
    \draw[blue,very thick] (-2.7,0.45)--(-0.7,0);
    \draw[blue,very thick] (-2.7,-0.45)--(-0.7,0);
    \draw[blue,very thick] (-0.7,0)
      .. controls (0.0,0.45) and (0.8,0.45) .. (1.5,0);
    \draw[blue,very thick] (-0.7,0)
      .. controls (0.0,-0.45) and (0.8,-0.45) .. (1.5,0);
    \draw[blue,very thick] (1.5,0)
      .. controls (1.1,0.85) and (1.9,0.85) .. (1.5,0);
    \fill[red] (-0.7,0) circle (2.2pt) node[below=3pt,text=black] {$z_1$};
    \fill[red] (1.5,0) circle (2.2pt) node[below=3pt,text=black] {$z_2$};
    \node[left=2pt] at (-2.7,0.45) {$x$};
    \node[left=2pt] at (-2.7,-0.45) {$y$};
  \end{tikzpicture}
  \diagramcaption{Dos líneas entre vértices y un lazo local.}
\end{minipage}
\end{center}
\begin{align}
  &\frac{(-i\lambda)^2}{4}
  \int d^4z_1\int d^4z_2\,
  \Delta_F(x-z_1)\Delta_F(y-z_1)
  \bigl[\Delta_F(z_1-z_2)\bigr]^2
  \Delta_F(z_2-z_2).
\end{align}

Diagrama \emph{sunset}:
\begin{center}
\begin{minipage}{\linewidth}
\centering
  \begin{tikzpicture}[scale=0.9,baseline=(current bounding box.center)]
    \draw[blue,very thick] (-2.7,0)--(-1.0,0);
    \draw[blue,very thick] (1.0,0)--(2.7,0);
    \draw[blue,very thick] (-1.0,0)
      .. controls (-0.45,0.75) and (0.45,0.75) .. (1.0,0);
    \draw[blue,very thick] (-1.0,0)--(1.0,0);
    \draw[blue,very thick] (-1.0,0)
      .. controls (-0.45,-0.75) and (0.45,-0.75) .. (1.0,0);
    \fill[red] (-1.0,0) circle (2.2pt) node[below=3pt,text=black] {$z_1$};
    \fill[red] (1.0,0) circle (2.2pt) node[below=3pt,text=black] {$z_2$};
    \node[below=3pt] at (-2.7,0) {$x$};
    \node[below=3pt] at (2.7,0) {$y$};
  \end{tikzpicture}
  \diagramcaption{Diagrama \emph{sunset}.}
\end{minipage}
\end{center}
\begin{equation}
  \frac{(-i\lambda)^2}{6}
  \int d^4z_1\int d^4z_2\,
  \Delta_F(x-z_1)
  \bigl[\Delta_F(z_1-z_2)\bigr]^3
  \Delta_F(z_2-y).
\end{equation}
Los denominadores $4$, $4$ y $6$ son los factores de simetría de estas topologías,
respectivamente. Se obtienen también contando las contracciones con los
puntos externos $x,y$ fijos: hay $288$, $288$ y $192$ emparejamientos,
respectivamente. Al dividir por el factor $2!(4!)^2=1152$ de la expansión
de Dyson, resultan $1/4$, $1/4$ y $1/6$.

\section{Reglas de Feynman en espacio de posiciones}

El cálculo anterior se resume en las reglas:
\begin{enumerate}[label=\textbf{\arabic*.},leftmargin=1.8em,itemsep=1.2em]
  \item \textbf{Propagador por cada línea.}
  Una línea que une los puntos $a$ y $b$ aporta $\Delta_F(a-b)$.
  \begin{center}
\begin{minipage}{\linewidth}
\centering
    \begin{tikzpicture}[baseline=(current bounding box.center)]
      \draw[blue,very thick] (-1.5,0)--(1.5,0);
      \fill[red] (-1.5,0) circle (2pt) node[below=3pt,text=black] {$a$};
      \fill[red] (1.5,0) circle (2pt) node[below=3pt,text=black] {$b$};
    \end{tikzpicture}
    \diagramcaption{Propagador entre dos puntos. La línea representa la contracción
    libre entre los campos situados en \(a\) y \(b\).}
  \end{minipage}
\end{center}
  \begin{equation}
    \Delta_F(a-b)=\mel{0}{T\{\phi_I(a)\phi_I(b)\}}{0}.
  \end{equation}

\item \textbf{Factor e integración por cada vértice.}
  Cada vértice de la interacción $\lambda\phi^4/4!$ aporta un factor
  $-i\lambda$, y su posición interna $z$ debe integrarse sobre todo el
  espacio-tiempo.
  \begin{center}
\begin{minipage}{\linewidth}
\centering
    \begin{tikzpicture}[baseline=(current bounding box.center)]
      \draw[blue,very thick] (-1.0,0)--(1.0,0);
      \draw[blue,very thick] (0,-0.65)--(0,0.65);
      \fill[red] (0,0) circle (2.5pt) node[above right=2pt,text=black] {$z$};
    \end{tikzpicture}
    \diagramcaption{Vértice de interacción. El punto \(z\) se integra sobre todo
    el espacio-tiempo y aporta el factor \((-i\lambda)\int d^4z\).}
  \end{minipage}
\end{center}
  \begin{equation}
    \text{Un vértice:}\quad(-i\lambda)\int d^4z.
  \end{equation}
  Con dos vértices internos $z_1,z_2$ aparecen dos integrales,
  \begin{equation}
    (-i\lambda)^2\int d^4z_1\int d^4z_2.
  \end{equation}

\item \textbf{Factor de simetría.}
  La contribución se divide por el factor de simetría $S$ del diagrama.
  \begin{center}
\begin{minipage}{\linewidth}
\centering
    \begin{tikzpicture}[baseline=(current bounding box.center)]
      \draw[blue,very thick] (-1.8,0)--(1.8,0);
      \draw[blue,very thick] (0,0) .. controls (-0.7,1.25) and (0.7,1.25) .. (0,0);
      \fill[red] (0,0) circle (2.5pt) node[below=3pt,text=black] {$z$};
      \node at (-1.8,-0.3) {$x$};
      \node at (1.8,-0.3) {$y$};
    \end{tikzpicture}
    \diagramcaption{Ejemplo de factor de simetría. La contracción interna en el
    mismo vértice produce el factor \(S=2\).}
  \end{minipage}
\end{center}
  \begin{equation}
    S=2,\qquad\frac{-i\lambda}{S}=\frac{-i\lambda}{2}.
  \end{equation}
  Por tanto, este diagrama \emph{tadpole} produce
  \begin{equation}
    -\frac{i\lambda}{2}\int d^4z\,
    \Delta_F(x-z)\Delta_F(y-z)\Delta_F(z-z).
  \end{equation}

\end{enumerate}

Para un diagrama con $V$ vértices, las tres reglas se resumen en
\begin{equation}
  \eqbox{\frac{1}{S}(-i\lambda)^V
  \int\prod_{j=1}^{V}d^4z_j
  \prod_{\text{líneas }(ab)}\Delta_F(a-b).}
\end{equation}

\subsection{Cancelación de burbujas de vacío}

Los diagramas desconectados sin patas externas se cancelan contra el
  denominador de Gell-Mann--Low. A primer orden, el numerador contiene tres
  clases de contribuciones: la propagación libre, esa misma línea acompañada
  por una burbuja desconectada y el diagrama \emph{tadpole} conectado.

\noindent\begin{minipage}{\linewidth}
  \begin{equation}
    \begin{tikzpicture}[scale=0.75,baseline=(current bounding box.center)]
      % Numerador.
      \node at (-6.1,2.0) {$N=$};
      \draw[blue,very thick] (-5.45,2.0)--(-4.25,2.0);
      \node[below=2pt] at (-5.45,2.0) {$x$};
      \node[below=2pt] at (-4.25,2.0) {$y$};
      \node at (-4.85,3.2) {\scriptsize libre};
      \node at (-3.88,2.0) {$+$};

      \draw[blue,very thick] (-3.5,2.0)--(-2.3,2.0);
      \node[below=2pt] at (-3.5,2.0) {$x$};
      \node[below=2pt] at (-2.3,2.0) {$y$};
      \draw[blue,very thick] (-1.62,2.0)
        .. controls (-2.05,2.55) and (-1.19,2.55) .. (-1.62,2.0);
      \draw[blue,very thick] (-1.62,2.0)
        .. controls (-2.05,1.45) and (-1.19,1.45) .. (-1.62,2.0);
      \fill[red] (-1.62,2.0) circle (2.2pt);
      \node at (-2.55,3.2) {\scriptsize desconectada};
      \node at (-0.98,2.0) {$+$};

      \draw[blue,very thick] (-0.58,2.0)--(0.92,2.0);
      \draw[blue,very thick] (0.17,2.0)
        .. controls (-0.28,3.05) and (0.62,3.05) .. (0.17,2.0);
      \fill[red] (0.17,2.0) circle (2.2pt);
      \node[below=2pt] at (-0.58,2.0) {$x$};
      \node[below=2pt] at (0.92,2.0) {$y$};
      \node at (0.17,3.2) {\scriptsize conectada};
      \node at (1.35,2.0) {$+\cdots$};

      % Denominador.
      \node at (-6.1,0.35) {$D=$};
      \node at (-5.35,0.35) {$1+$};
      \draw[blue,very thick] (-4.35,0.35)
        .. controls (-4.8,0.9) and (-3.9,0.9) .. (-4.35,0.35);
      \draw[blue,very thick] (-4.35,0.35)
        .. controls (-4.8,-0.2) and (-3.9,-0.2) .. (-4.35,0.35);
      \fill[red] (-4.35,0.35) circle (2.2pt);
      \node at (-3.72,0.35) {$+\cdots$};
      \node at (-4.35,1.15) {\scriptsize vacío};

      % Cociente: sólo sobreviven las piezas conectadas.
      \node at (-6.0,-1.25) {$\dfrac ND=$};
      \draw[blue,very thick] (-5.1,-1.25)--(-3.9,-1.25);
      \node[below=2pt] at (-5.1,-1.25) {$x$};
      \node[below=2pt] at (-3.9,-1.25) {$y$};
      \node at (-3.5,-1.25) {$+$};
      \draw[blue,very thick] (-3.05,-1.25)--(-1.55,-1.25);
      \draw[blue,very thick] (-2.3,-1.25)
        .. controls (-2.75,-0.2) and (-1.85,-0.2) .. (-2.3,-1.25);
      \fill[red] (-2.3,-1.25) circle (2.2pt);
      \node[below=2pt] at (-3.05,-1.25) {$x$};
      \node[below=2pt] at (-1.55,-1.25) {$y$};
      \node at (-1.1,-1.25) {$+\cdots$};
      \node[blue] at (-3.45,-2.4) {\small sólo permanecen diagramas conectados};
    \end{tikzpicture}
  \end{equation}
  \begin{center}
\begin{minipage}{\linewidth}
\centering
    \diagramcaption{Cancelación de burbujas de vacío en el cociente \(N/D\).}
  \end{minipage}
\end{center}
  \end{minipage}

En $N$, la contribución de la burbuja desconectada multiplica al
propagador externo; en $D$ aparece la misma contribución de vacío.
Es este factor común, con sus coeficientes correspondientes, el que se
cancela al formar el cociente.

Para obtener el correlador normalizado, se suman los diagramas conectados
al par externo $x,y$ al orden deseado, con esos puntos fijos. Las burbujas
de vacío ya se han cancelado según \eqref{eq:VacuumDiagramCancellation};
los lazos unidos a los puntos externos permanecen. El factor $1/S$ incluye
las simetrías del diagrama y no debe aplicarse de nuevo si el coeficiente
ya se obtuvo contando las contracciones de Wick.

\end{document}
